\section{Reconstruction error: what each measured quantity changes}
\label{sec:measurement}
First keep the coordinate frame correct, and perturb only the reconstruction
inputs. Define used-minus-true errors by
\begin{equation}
 \hat R_s=\Rs+\resistanceerror,\quad \hat\voltage=\voltage+\vect\varepsilon_u,
 \quad\hat\current=\current+\vect\varepsilon_i,\quad
 \hat\omega_e=\omegae+\varepsilon_\omega.
 \label{eq:error-definitions}
\end{equation}
Substitute these values into \cref{eq:reconstruction}. Before linearizing,
the difference at the same physical operating point is exactly
\begin{equation}
 \hat\flux^{\rm rec}-\flux=
 \frac{-\rotationgenerator(\vect\varepsilon_u-\Rs\vect\varepsilon_i
       -\resistanceerror\current-\resistanceerror\vect\varepsilon_i)
       -\varepsilon_\omega\flux}{\omegae+\varepsilon_\omega}.
 \label{eq:exact-input-error}
\end{equation}
Its differential, neglecting products of small input errors and assuming
$|\varepsilon_\omega|\ll|\omegae|$, is
\begin{equation}
 \boxed{\delta\flux_{\rm rec}=
 -\frac{\rotationgenerator\delta\voltage}{\omegae}
 +\frac{\Rs\rotationgenerator\delta\current}{\omegae}
 +\frac{\rotationgenerator\current}{\omegae}\delta R_s
 -\frac{\flux}{\omegae}\delta\omega_e.}\label{eq:input-sensitivity}
\end{equation}
Taking components checks every sign:
\begin{align}
 \delta\psid&=\frac{\delta u_q-\Rs\delta i_q-i_q\delta R_s
                  -\psid\delta\omega_e}{\omegae},\\
 \delta\psiq&=\frac{-\delta u_d+\Rs\delta i_d+i_d\delta R_s
                  -\psiq\delta\omega_e}{\omegae}.
\end{align}
The rotation converts q-voltage error into d-flux error and d-voltage error
into negative q-flux error. Differentiating the reciprocal speed supplies the
negative speed term. All numerators have units of volts, since
$\psi\delta\omega_e$ is a voltage; division by angular speed gives webers.

\subsection{A current error also changes the map label}
The preceding differential concerns a value at one physical point.
When measurements are stored at $\hat\current$, the true magnetic value was
evaluated at $\hat\current-\vect\varepsilon_i$. Comparing maps at the same
numerical label introduces an additional displacement term:
\begin{equation}
 \delta\flux_{i,\rm map}=
 \left(\frac{\Rs}{\omegae}\rotationgenerator-\Ldiff\right)\vect\varepsilon_i.
 \label{eq:current-relabel}
\end{equation}
Ignoring this term could misinterpret a shifted magnetic curve as a pure
voltage-reconstruction error. A coherent frame rotation will be treated
separately, because it rotates voltage and flux as well as current.

\subsection{Resistance mismatch}
With perfect signals and $\resistanceerror=\hat R_s-\Rs$, the resistance contribution
is exact, not merely first order:
\begin{equation}
 \boxed{\delta\psi_{d,R}=-\frac{i_q}{\omegae}\resistanceerror,\qquad
        \delta\psi_{q,R}=+\frac{i_d}{\omegae}\resistanceerror.}\label{eq:R-error}
\end{equation}
At fixed $i_d$ and paired equal speed, the first is odd and the second even
in $i_q$. At $i_d=0$ the q contribution is zero. Overestimating resistance at
positive speed therefore adds a negative d-flux slope. A correction inferred
as $\resistanceerror>0$ must be \emph{subtracted} from the used resistance.
\begin{figure}[tb]
\centering
\begin{tikzpicture}
\begin{axis}[diagnostic axis,title={Subtracting too much resistive voltage},
 ylabel={$\hat\psi_d/\psi_*$},xmin=-1,xmax=1,ymin=.975,ymax=1.055]
 \addplot[reference quantity,strong line,domain=-1:1]{1+.02*x^2};
 \addlegendentry{Ideal}
 \addplot[estimated quantity,standard line,domain=-1:1]{1+.02*x^2-.025*x};
 \addlegendentry{$\varepsilon_R>0$}
\end{axis}
\begin{axis}[diagnostic axis,at={(.5\textwidth,0)},anchor=south west,
 title={Forbidden channels at $i_d=0$},ylabel={$r/\psi_*$},
 xmin=-1,xmax=1,ymin=-.03,ymax=.03]
 \addplot[estimated quantity,strong line,domain=-1:1]{-.025*x};
 \addlegendentry{$r_d$}
 \addplot[derived quantity,alternative pattern,standard line,domain=-1:1]{0};
 \addlegendentry{$r_q=0$}
\end{axis}
\end{tikzpicture}
\caption{Resistance-only reconstruction error at $\normalizedspeed=1$, $x=0$,
and $\varepsilon_R/R_*=0.025$. The ideal nonlinear d flux remains even;
incorrect resistive subtraction adds a linear odd contribution.}
\label{fig:resistance}
\end{figure}

This tilt originates in subtracting too much resistive voltage; no change to
the magnetic constitutive law is involved. Resistance and inverter uncertainty
are recognized obstacles to flux-map identification \cite{liu2018}.

At a fixed nonzero speed, the same resistance error also violates integrability:
\begin{equation}
 r_{\rm curl,R}=\partial_{i_d}\delta\psi_{q,R}
               -\partial_{i_q}\delta\psi_{d,R}=2\resistanceerror/\omegae.
 \label{eq:resistance-curl}
\end{equation}
Thus resistance error can violate both parity and reciprocity, whereas the
pure coherent angle rotation derived next preserves reciprocity. Derivative
residuals amplify noise and require fixed speed across the current-plane
derivative. If speed varies with current, derivatives of $1/\omegae$ enter
and this simple curl expression no longer applies.
